% SOURCE_AUTHORITY: sga5_fr_workpass.tex, lines 8258--8286 (LF-normalized unit).
% SOURCE_SHA256: 791F4EFFC5E02832D5D77ED03518C8156D6F07E4C8238B03545DB93D883FBB28
\subsubsection*{1.1.4}
Seja \(F=(F_n)_{n\in\mathbb N}\) um feixe \(\ell\)-ádico construtível. Como o anel \(\mathbb Z_\ell\) atua de forma evidente sobre cada feixe \(F_n\), também atua sobre \(F\). Além disso, se \(F\) e \(G\) são dois feixes \(\ell\)-ádicos construtíveis e
\[
u:F\longrightarrow G
\]
é um morfismo, \(u\) é compatível com as ações de \(\mathbb Z_\ell\) sobre \(F\) e \(G\), respectivamente. Isso mostra que \(\ell\text{-}\adc(X)\) é canonicamente dotada de uma estrutura de \(\mathbb Z_\ell\)-categoria abeliana.

Isto justifica a seguinte terminologia:

\begin{statement}{Terminologia}
Os feixes \(\ell\)-ádicos construtíveis serão chamados \(\mathbb Z_\ell\)-feixes construtíveis.
\end{statement}

É conveniente não confundir esta noção com a de \((\mathbb Z_\ell)_X\)-módulo, que não será usada no que segue.

A categoria \(\ell\text{-}\adc(X)\) será igualmente indicada por
\[
\mathbb Z_\ell\text{-}\fc(X).
\]

\subsubsection*{1.2. \(\mathbb Z_\ell\)-feixes constantes torcidos construtíveis}

\begin{statement}{Definição 1.2.1}
Um \(\mathbb Z_\ell\)-feixe construtível
\[
F=(F_n)_{n\in\mathbb N}
\]
é denominada constante torcida se e somente se os feixes \(F_n\) forem localmente constantes sobre \(X_{\mathrm{\acute et}}\).
\end{statement}
